\subsection{Оптимизация производительности}

При проектировании системы учитывалось, что объём данных может достигать сотен тысяч животных и десятков миллионов записей. Для обеспечения приемлемого времени отклика при работе с набором данных FelBenitez~\cite{FelBenitez}, содержащим более 210\,000 животных и 12\,600\,000 записей надоев, был проведён комплексный анализ узких мест и реализован ряд оптимизаций на всех уровнях архитектуры.

\subsubsection{Анализ узких мест}

Профилирование запросов к базе данных с помощью \texttt{EXPLAIN ANALYZE} выявило три основные категории проблем:

\begin{enumerate}
    \item \textbf{Коррелированные подзапросы} --- каждый предиктивный эндпоинт выполнял от 10 до 17 \texttt{LATERAL JOIN} на каждое животное. Для 210\,000 животных и 15 подзапросов это составляет свыше 3\,000\,000 executions подзапросов за один HTTP-запрос.
    \item \textbf{Повторная загрузка моделей} --- функция \texttt{joblib.load()} вызывалась при каждом предикте, выполняя чтение с диска и десериализацию объекта XGBoost.
    \item \textbf{Отсутствие кеширования} --- результаты аналитических запросов и ML-предиктов не кешировались; повторные запросы выполняли полный цикл расчётов.
\end{enumerate}

В таблице~\ref{tab:bottlenecks} приведены типичные времена ответа до оптимизации для ключевых эндпоинтов.

\begin{table}[htbp]
\caption{Время ответа ключевых эндпоинтов до оптимизации}
\label{tab:bottlenecks}
\centering
\begin{tabular}{|l|c|}
\hline
\textbf{Эндпоинт} & \textbf{Время ответа, с} \\
\hline
Дашборд (\texttt{/analytics/dashboard}) & 5--10 \\
\hline
Предикт мастита (одно животное) & 30--60 \\
\hline
Предикт мастита (всё стадо) & 180--300 \\
\hline
Временная шкала здоровья (90 дней) & 10--30 \\
\hline
Загрузка модели с диска & 0{,}05--0{,}2 \\
\hline
\end{tabular}
\end{table}

\subsubsection{Оптимизация базы данных}

Для ускорения аналитических запросов к TimescaleDB реализован комплекс мер: композитные индексы, автоматическое сжатие исторических данных и материализованные агрегаты.

\textbf{Композитные индексы.} Большинство коррелированных подзапросов используют паттерн \texttt{WHERE animal\_id = \$1 ORDER BY date DESC LIMIT 1}. Для данного паттерна созданы составные индексы, позволяющие выполнить запрос одним сканированием индекса:

\begin{verbatim}
CREATE INDEX idx_calvings_animal_date_desc
    ON calvings (animal_id, calving_date DESC);
CREATE INDEX idx_inseminations_animal_date_desc
    ON inseminations (animal_id, insemination_date DESC);
CREATE INDEX idx_vet_records_animal_type_confirmed_date
    ON vet_records (animal_id, record_type, confirmed, event_date);
\end{verbatim}

\textbf{Сжатие TimescaleDB.} Для таблиц с историческими данными включено автоматическое сжатие chunks старше 6 месяцев с сегментацией по \texttt{animal\_id}. Это обеспечивает до 550-кратного сжатия (например, таблица \texttt{milk\_day\_productions} --- с 1633~МБ до 2{,}96~МБ).

\textbf{Непрерывные агрегаты.} Созданы 5 материализованных представлений, предварительно вычисляющих суточные и недельные агрегаты:

\begin{verbatim}
CREATE MATERIALIZED VIEW milk_daily_summary
WITH (timescaledb.continuous) AS
SELECT time_bucket('1 day', date) AS bucket,
       animal_id,
       AVG(milk_amount) AS avg_milk,
       SUM(milk_amount)  AS total_milk
FROM milk_day_productions
GROUP BY bucket, animal_id;
\end{verbatim}

Агрегаты обновляются автоматически каждый час. Использование предвычисленных агрегатов вместо сканирования сырых данных сокращает объём обрабатываемых строк в 100--1000 раз.

\subsubsection{Кеширование моделей машинного обучения}

Для устранения повторной загрузки моделей реализован модуль кеширования, загружающий модели в оперативную память при первом обращении и проверяющий время модификации файла (\texttt{mtime}) для автоматической перезагрузки при переобучении.

\begin{verbatim}
_model_cache: dict[str, tuple[dict, float, str]] = {}

def get_model(name: str, path: str) -> dict | None:
    mtime = os.path.getmtime(path)
    if name in _model_cache:
        cached_data, cached_mtime, _ = _model_cache[name]
        if cached_mtime == mtime:
            return cached_data  # cache hit
    data = joblib.load(path)    # cache miss
    _model_cache[name] = (data, mtime, path)
    return data
\end{verbatim}

Аналогичный подход применён для ONNX Runtime сессий: \texttt{InferenceSession} создаётся один раз и переиспользуется. Это сокращает время инференса с 50--200~мс (загрузка + десериализация) до менее 5~мс.

Для вектора признаков (\texttt{DataFrame}) реализован TTL-кеш на 60~секунд, предотвращающий повторное выполнение SQL-запроса с 15 коррелированными подзапросами при нескольких предиктах одного типа в течение минуты.

\subsubsection{Интеграция колоночной базы данных ClickHouse}

Для аналитических запросов типа \texttt{GROUP BY animal\_id} с агрегацией по 12~миллионам строк интегрирована колоночная база данных ClickHouse. Колоночное хранение позволяет выполнять агрегатные запросы в 50--100 раз быстрее PostgreSQL, так как считываются только необходимые столбцы, а не вся строка.

\textbf{Архитектура синхронизации.} PostgreSQL остаётся источником истины для транзакционных операций (CRUD, синхронизация с Lely). ClickHouse выступает аналитической репликой только для чтения. Синхронизация выполняется каждые 5~минут батчами по 50\,000 строк через HTTP API ClickHouse.

\begin{figure}[htbp]
\centering
\includegraphics[width=0.85\textwidth,height=0.5\textheight,keepaspectratio]{clickhouse-arch.png}
\caption{Архитектура синхронизации PostgreSQL $\to$ ClickHouse}
\label{fig:clickhouse-arch}
\end{figure}

В таблице~\ref{tab:ch-schema} приведена схема таблиц в ClickHouse.

\begin{table}[htbp]
\caption{Аналитические таблицы ClickHouse}
\label{tab:ch-schema}
\centering
\begin{tabular}{|l|l|c|}
\hline
\textbf{Таблица} & \textbf{Ключ сортировки} & \textbf{Партиция} \\
\hline
\texttt{milk\_day\_productions} & \texttt{(animal\_id, date)} & \texttt{YYYYMM} \\
\hline
\texttt{milk\_quality} & \texttt{(animal\_id, date)} & \texttt{YYYYMM} \\
\hline
\texttt{ruminations} & \texttt{(animal\_id, date)} & \texttt{YYYYMM} \\
\hline
\texttt{feed\_day\_amounts} & \texttt{(animal\_id, date)} & \texttt{YYYYMM} \\
\hline
\texttt{activities} & \texttt{(animal\_id, date)} & \texttt{YYYYMM} \\
\hline
\end{tabular}
\end{table}

Пример запроса к ClickHouse, заменяющего 15 коррелированных подзапросов одним проходом:

\begin{verbatim}
SELECT animal_id,
    avgIf(milk_amount, date >= today() - 7) as avg_milk_7d,
    avgIf(scc, date >= today() - 30)       as avg_scc_30d,
    avgIf(rumination_minutes,
          date >= today() - 7)              as avg_rumination_7d
FROM milk_day_productions m
LEFT JOIN milk_quality q USING (animal_id, date)
LEFT JOIN ruminations r  USING (animal_id, date)
GROUP BY animal_id
\end{verbatim}

\subsubsection{Оптимизация серверной части}

\textbf{Кеширование в Redis.} Реализован модуль \texttt{MlCache} на основе Redis с настраиваемым TTL. Результаты дашборда кешируются на 1~час, результаты ML-предиктов --- на 1~час с ключом, включающим модель и текущий час. При промахе кеша запрос выполняется обычным образом, после чего результат записывается в Redis.

\textbf{Оконные функции вместо коррелированных подзапросов.} Эндпоинт \texttt{health\_timeline} ранее выполнял \texttt{generate\_series} с 4 коррелированными подзапросами на каждый из 365 дней ($365 \times 4 = 1460$ подзапросов). Запрос переписан с использованием оконных функций:

\begin{verbatim}
SELECT date,
  AVG(milk_amount) OVER (
    ORDER BY date ROWS BETWEEN 2 PRECEDING
                        AND CURRENT ROW) as short_milk,
  AVG(milk_amount) OVER (
    ORDER BY date ROWS BETWEEN 29 PRECEDING
                        AND CURRENT ROW) as long_milk,
  STDDEV(milk_amount) OVER (
    ORDER BY date ROWS BETWEEN 29 PRECEDING
                        AND CURRENT ROW) as milk_std
FROM milk_day_productions
WHERE animal_id = $1
ORDER BY date
\end{verbatim}

Такой запрос выполняет один проход по данным вместо 1460 коррелированных подзапросов, что сокращает время выполнения в 10--30 раз.

\textbf{Фоновый мониторинг дрейфа.} Вычисление статистик дрейфа данных (среднее, стандартное отклонение, квантили по 210\,000 строк и 20 признакам) перенесено из синхронного кода в фоновые задачи FastAPI (\texttt{BackgroundTasks}), что устраняет задержку при отправке ответа клиенту.

\subsubsection{Оптимизация клиентского приложения}

Для таблиц с большим количеством строк на странице аналитики применены два подхода:

\begin{enumerate}
    \item \textbf{Реактивные вычисления} --- сортировка и фильтрация данных вынесены в декларативные выражения \texttt{\$derived}, которые вычисляются только при изменении исходных данных, а не при каждом рендеринге компонента.
    \item \textbf{Виртуальный скроллинг} --- библиотека \texttt{svelte-tiny-virtual-list} обеспечивает рендеринг только видимых строк таблицы, что сокращает количество DOM-узлов с тысяч до 20--30 при любой прокрутке.
\end{enumerate}

\subsubsection{Результаты оптимизации}

В таблице~\ref{tab:perf-results} приведено сравнение времени ответа ключевых эндпоинтов до и после оптимизации на наборе данных 10\,435 животных (измерено с помощью \texttt{curl -w}).

\begin{table}[htbp]
\caption{Сравнение времени ответа до и после оптимизации (10\,435 животных)}
\label{tab:perf-results}
\centering
\begin{tabular}{|l|c|c|c|}
\hline
\textbf{Эндпоинт} & \textbf{До, с} & \textbf{После, с} & \textbf{Ускорение} \\
\hline
Дашборд & 5--10 & 0{,}13 & $38--77\times$ \\
\hline
ML-предикт мастита (1 животное) & 5{,}0 & 0{,}04 & $125\times$ \\
\hline
ML-предикт мастита (всё стадо) & 5{,}0 & 4{,}6 & $1{,}1\times$ \\
\hline
ML-предикт эструса (1 животное) & 5{,}0 & 0{,}03 & $167\times$ \\
\hline
ML-предикт кетоза (1 животное) & 1{,}2 & 0{,}02 & $60\times$ \\
\hline
ML-предикт кормления (1 животное) & 0{,}34 & 0{,}01 & $34\times$ \\
\hline
ML-предикт выбраковки (1 животное) & 2{,}3 & 0{,}02 & $115\times$ \\
\hline
Временная шкала здоровья (90 дней) & 10--30 & 0{,}008 & $1250--3750\times$ \\
\hline
\end{tabular}
\end{table}

Основное ускорение при предикте для одного животного достигнуто за счёт фильтрации по \texttt{animal\_id} на уровне SQL-запроса к загрузчику данных (\texttt{WHERE a.id = :aid}), вместо загрузки признаков для всех 10\,435 животных и последующей фильтрации в памяти. Для эндпоинтов, запрашивающих предикт для всего стада, время ответа остаётся значительным (4--5~с), что обусловлено объёмом вычислений для 10\,000+ животных; для таких запросов рекомендуется асинхронный режим с кешированием результата в Redis.

% TODO: uncomment when benchmark data is available
% \begin{figure}[htbp]
% \centering
% \includegraphics[width=0.9\textwidth,height=0.5\textheight,keepaspectratio]{perf-benchmark.png}
% \caption{Результаты нагрузочного тестирования до и после оптимизации}
% \label{fig:perf-benchmark}
% \end{figure}
